% clase article a una cara, con letra de 12pt, en papel A4
% y con p�gina separada para el t�tulo
\documentclass[oneside,12pt,a4paper,titlepage]{article}

% Carga de paquetes necesarios. "customarticle" es un paquete personalizado
\usepackage[spanish,english]{babel} % Internacionalizaci�n de LaTeX
\usepackage[ansinew]{inputenx}% C�digo de caracteres de teclado
\usepackage[T1]{fontenc}      % Codificaci�n de fundiciones
\usepackage{times}            % Familia base de fundiciones
\usepackage{graphicx}
\usepackage{pifont}
\usepackage[usenames,dvipsnames]{color}
\usepackage{amsmath}
\usepackage{hyperref}
\usepackage{colortbl}
\usepackage{makeidx}
\usepackage{indentfirst}
\hypersetup{bookmarksopen,bookmarksopenlevel=3,linktocpage,colorlinks,urlcolor=blue,citecolor=blue,linkcolor=black,filecolor=blue,pdfnewwindow, 	pdftitle={DESGLOSA: Un sistema de visualizaci�n 3D para dar soporte al Desarrollo Global de Software}, pdfsubject={Anteproyecto} pdfauthor={Jose Domingo L�pez L�pez}}
\usepackage[spanish,caption,counter,title,fancy]{customarticle}
\usepackage[fixlanguage]{babelbib}
\usepackage{url}
%% Define a new 'leo' style for the package that will use a smaller font.
\makeatletter
\def\url@leostyle{%
\@ifundefined{selectfont}{\def\UrlFont{\sf}}{\def\UrlFont{\footnotesize\ttfamily}}}
\makeatother

\makeindex

% Macro para insertar una imagen
%       Uso: \imagen{nombreFichero}{Factor escala}{Caption (leyenda)}{Label (identificador para referenciarla)}
\def\imagen#1#2#3#4{
 \begin{figure}[h]
 \begin{center}
   \includegraphics[keepaspectratio,scale=#2]{#1}
 \caption {#3}
 \label{#4}
 \end{center}
 \end{figure}
}

% Macro para insertar una imagen con borde en el marco
%       Uso: \imagenBorde{nombreFichero}{Factor escala}{Caption (leyenda)}{Label (identificador para referenciarla)}
\def\imagenBorde#1#2#3#4{
\setlength\fboxsep{0pt}
\setlength\fboxrule{0.5pt}
 \begin{figure}[h]
 \begin{center}
  \fbox{\includegraphics[keepaspectratio,scale=#2]{#1}}
 \caption{#3}
 \label{#4}
 \end{center}
\end{figure}}

% Definici�n de una lista personalizada 
\newenvironment{myitemize}%
{\begin{list}{\textbullet}%
{\settowidth{\labelwidth}{\textbullet}
\addtolength{\topsep}{-1.0ex}
\addtolength{\itemsep}{\dimexpr 0ex plus 0.15ex minus 0.25ex}
\setlength{\parsep}{0.0ex}
\setlength{\partopsep}{1.75ex}
\setlength{\parskip}{0.0ex}
}}%
{\end{list}}

\begin{document}
\selectlanguage{spanish}
% En las p�ginas de portada e �ndices, no hay encabezado ni pie de p�gina
\pagestyle{empty}
\begin{titlepage}
	\begin{center}
		\includegraphics[scale=0.25, keepaspectratio]{./images/esi_bw.png} \\
		\vspace{20mm}
  	{\large \textbf{UNIVERSIDAD DE CASTILLA-LA MANCHA}} \\
  	\bigskip
  	{\large \textbf{ESCUELA SUPERIOR DE INFORM�TICA}} \\
  	\vspace{25mm}
  	{\large \textbf{INGENIER�A}}\\
  	\bigskip
  	{\large \textbf{EN INFORM�TICA}}\\
  	\vspace{25mm}
  	{\large \textbf{ANTEPROYECTO}}\\
  	\vspace{15mm}
  	{Cuadro de Mando para la Visualizaci�n de la Calidad del Software, Integrado con un Entorno de Medici�n}\\
  	\vspace{15mm}
  	{Jose Domingo L�pez L�pez}
  	\vspace{20mm}
  	\begin{flushright}
  		\textbf{Septiembre, 2010}
  	\end{flushright}
	\end{center}
\end{titlepage}
% -------------
%% Derechos de autor
%\parskip=10pt
%\mbox{}
%\vfill{}
%\begin{small}
%	\noindent DESGLOSA\\
%	\noindent Un sistema de visualizaci�n 3D para dar soporte al Desarrollo Global de Software
%
%	\noindent\copyright~ 2011 Jose Domingo L�pez L�pez (Alumno del PFC)\\
%	\phantom{\copyright~ 2011} M\textordfeminine{} �ngeles Moraga de la Rubia (Director del PFC)
%
%	\noindent En virtud del Art�culo 7 del Real Decreto Legislativo 1/1996, de 12 de abril por el que se aprueba el Texto Refundido de la Ley de Propiedad Intelectual, modificado por la Ley 23/2006 de 7 de julio, este PFC se considera una obra en colaboraci�n entre las diferentes partes. Por tanto la propiedad intelectual de este PFC ser� compartida con iguales porcentajes entre el alumno y el director del PFC.
%
%	\noindent Este documento ha sido compuesto con \LaTeX{} y sus im�genes generadas con Inkscape y GIMP.
%\end{small}
%\clearpage
% -------------
% En los �ndices numeramos las p�ginas en n�meros romanos
\pagestyle{plain}
\pagenumbering{Roman}
% �ndice de secciones
\tableofcontents
\clearpage

\parskip=10pt
%\onehalfspacing   % Espaciado 1,5 l�neas (facilitado por setspace)
\pagenumbering{arabic}
\pagestyle{fancy}
% Aqui se incluyen los archivos .tex que forman el documento
\section{Introducci�n}

A�os atr�s, las organizaciones y empresas se preocupaban de la calidad de los procesos que se segu�an para desarrollar software bas�ndose en la premisa de que \textit{``con un proceso de calidad se obtiene un producto de calidad''} \cite{evaluacion-calidad}, pero dados los problemas que han surgido en sistemas ya implantados y en producci�n, y las p�rdidas que �stos han ocasionado, ha sido necesario actualizar la percepci�n que se tiene acerca de la calidad del software. Por ello, las organizaciones comenzaron a preocuparse por la calidad del producto software desde el punto de vista de los clientes y usuarios finales, abarcando un mayor n�mero de \textit{stakeholders}, con el fin de aumentar el grado de satisfacci�n y confianza por parte de los mismos.

Es por esto por lo que se tratado de asegurar la funcionalidad y usabilidad de los sistemas profundizando en �rea de \textit{testing}, pero se han quedado sin cubrir aspectos de seguridad y mantenibilidad. Y es que un producto software no es �nicamente el c�digo que se ejecuta sobre una m�quina, si no que lo son tambi�n los artefactos que se han ido produciendo a lo largo de todo el ciclo de vida: modelos, diagramas, documentos, etc.

Para paliar algunos de los problemas expuestos anteriormente, se comenzaron a utilizar dos est�ndares:
\begin{myitemize}
	\item \textbf{ISO/IEC 9126} \cite{iso9126}, que define el modelo de calidad de un producto, incluyendo las caracter�sticas, subcaracter�sticas y m�tricas (de calidad interna, calidad externa y calidad en uso) que han de tenerse en cuenta, pero sin proporcionar una metodolog�a para su evaluaci�n. Las caracter�sticas que propone son funcionalidad, fiabilidad, usabilidad, eficiencia, mantenibilidad y portabilidad.
	\item \textbf{ISO/IEC 14598} \cite{iso14598}, que proporciona gu�as y requisitos para el proceso de evaluaci�n, cubriendo as� los aspectos de los que carece la norma ISO/IEC 9126.
\end{myitemize}

Pero entre ambos est�ndares existen distintas incompatibilidades e inconsistencias -tales como diferencias en el vocabulario, t�rminos y conceptos, la necesidad de a�adir recomendaciones, metodolog�as, gu�as y nuevas necesidades para especificar las dimensiones de la calidad del software, entre otros- que motivan la creaci�n de la familia \textbf{ISO/IEC 25000}, denominada SQuaRe (\underline{S}oftware product \underline{Qua}lity \underline{R}equirements and \underline{E}valuation).

La norma ISO/IEC 25000 establece criterios para la especificaci�n de requisitos de calidad de productos software y su evaluaci�n, un modelo de calidad y m�tricas para los distintos atributos de dicho modelo. N�tese que no es necesario alcanzar una calidad perfecta o absoluta, sino la necesaria y suficiente para cada contexto de uso y en funci�n de las necesidades de los \textit{stakeholders}, que son todos las entidades involucradas en el producto: clientes, usuarios finales, factor�as y empresas de desarrollo, etc�tera.
% TODO En este p�rrafo sacar lo del post de mois�s
% http://www.iso15504.es/index.php/component/content/article/82-blog/165-la-familia-de-normas-iso-25000-y-su-relacion-con-iso-15504-spice.html

La aparici�n de esta norma, cuyo objetivo es reemplazar a las normas ISO/IEC 9126 e ISO/IEC 14598, junto con las nuevas necesidades que surgen en las empresas para controlar y evaluar la calidad de los desarrollos inform�ticos, motivan la aparici�n del proyecto \textbf{MEDUSAS} (\underline{M}ejora y \underline{E}valuaci�n del \underline{D}ise�o, \underline{U}sabilidad, \underline{S}eguridad y \underline{M}antenibilidad del \underline{S}oftware). MEDUSAS (Fig. \ref{fig:arch-medusas}) es un entorno metodol�gico e instrumental basado en la ISO/IEC 25000 para satisfacer dichas necesidades, y evaluar la calidad del software teniendo en cuenta todos y cada uno de sus artefactos: c�digo fuente, dise�o, diagramas, modelos, etc.

\imagen{./images/comp_medusas.png}{0.4}{Arquitectura del proyecto MEDUSAS}{fig:arch-medusas}

Aqu� es donde surge la motivaci�n del Proyecto Fin de Carrera (en adelante PFC), que consistir� en realizar un cuadro de mando para la visualizaci�n de la calidad del software, integrado con el entorno MEDUSAS. Como resultado de la integraci�n, se dispondr� de un sistema completo que permitir� automatizar la adquisici�n, presentaci�n y an�lisis \cite{gil95} de las m�tricas que definen la calidad del producto software.

\clearpage
\section{Fundamentos te�ricos}

Para la realizaci�n de este PFC es necesario definir y estudiar tres conceptos fundamentales:
\begin{myitemize}
	\item Qu� es el Desarrollo Global de Software (GSD).
	\item Qu� son las m�tricas y para qu� se utilizan.
	\item C�mo visualizar m�tricas de un modo intuitivo, efectivo y eficaz.
\end{myitemize}


\subsection{Desarrollo Global de Software}
Sahay Sundeep (2003) define el GSD como \textit{``una forma coordinada de desarrollo de software realizado en lugares geogr�ficamente distantes atravesando fronteras nacionales, implicando interacci�n s�ncrona (en tiempo real) y as�ncrona''}. Por otro lado, Parastoo Mohagheghi (2004) lo define como el \textit{``modo de desarrollo de software que permite que los proyectos, productos o sistemas atraviesen las fronteras de las compa��as y pa�ses, por medio de la coordinaci�n de equipos distribuidos con distintas localizaciones geogr�ficas y husos horarios''}. Tomando como base ambas definiciones, se puede generalizar y definir el GSD como \textit{un conjunto de actividades que permite desarrollar software mediante equipos geogr�ficamente distantes, que pueden pertenecer a diferentes organizaciones, y que comparten flujos de trabajo y entregables}. Esta �ltima parte cobra fuerza si se reflexiona acerca de que el ciclo de desarrollo del software se compone de una serie de fases relacionadas entre s� (an�lisis, dise�o, implementaci�n, pruebas, despliegue, mantenimiento y soporte) y que un producto software no es �nicamente el c�digo fuente que se ejecuta en una o varias m�quinas, si no que lo conforman todos los artefactos que se han ido generando a lo largo de su ciclo de vida (diagramas, modelos, documentaci�n, etc�tera).

Por tanto, tal y como se ha definido el GSD, se infiere que el desarrollo distribuido afecta tanto al proceso de desarrollo del software como a la gesti�n de proyectos, y es susceptible de problemas tales como las diferencias horarias, culturales o ling��sticas, coordinaci�n, motivaci�n, etc�tera \cite{herb-mockus-2001}. Pero tambi�n proporciona una serie de ventajas muy importantes que prometen, entre otros, beneficios econ�micos \cite{herb-moitra-2001}. A continuaci�n se enumeran algunos de ellos:
\begin{myitemize}
	\item \textbf{Reducci�n de costes:} gracias a la subcontrataci�n de empresas o equipos de desarrollo ubicados en localidades o pa�ses en los que la mano de obra es m�s econ�mica.
	\item \textbf{Tiempos de lanzamiento al mercado m�s cortos:} ya que se pueden subcontratar varios equipos de desarrollo para desempe�ar distintas actividades de un modo paralelo o concurrente.
	\item \textbf{Mejora de la calidad:} se debe a que determinadas funcionalidades pueden ser desarrolladas por expertos ajenos a la empresa que desarrolla el software.
	\item \textbf{Escalabilidad:} hace posible afrontar la alza o baja, tanto temporal como permanente, de demanda de productos
\end{myitemize}


\subsection{M�tricas}



\subsection{Visualizaci�n de m�tricas}

Los sistemas de visualizaci�n deben proporcionar mecanismos que permitan que el an�lisis de datos sean tareas sean r�pidas y precisas. Para lograrlo, se hace uso de una cualidad del sistema visual humano denominado \textit{proceso pre-atentivo} \cite{hea92}. Los procesos pre-atentivos son procesos que ``saltan a la vista'', siendo detectados inmediatamente y evitando que el usuario deba centrar su atenci�n en una determinada regi�n de una imagen para captar la presencia o ausencia de una caracter�stica determinada \cite{ducasse10}.

Existen cuatro categor�as b�sicas de caracter�sticas que se procesan de forma pre-atentiva \cite{ware99}: color, forma, movimiento y localizaci�n espacial. De cada una de estas caracter�sticas se puede obtener un gran conjunto de subcaracter�sticas como curvatura, tama�o, orientaci�n, longitud, etc�tera.

Por otro lado, gracias a las \textit{vistas polim�tricas} \cite{lanza03} es posible representar una m�trica de un modo visual. El m�todo consiste en establecer un conjunto de asociaciones entre cada una de las dimensiones de la m�trica con cada una de las subcaracter�sticas de un modelo gr�fico. De este modo, los modelos tomar�n una determinada forma, orientaci�n o color en funci�n de los valores de las m�tricas, permitiendo as� hacer uso de los beneficios de los procesos pre-atentivos.
\clearpage
\section{Objetivos}

Los objetivos que se esperan lograr con la realizaci�n de este PFC se han dividido en dos categor�as. En la primera de ellas, como objetivo principal, se introducir� el producto final que se espera obtener y, por �ltimo, como objetivos secundarios, se expondr�n los objetivos que se ir�n alcanzando en las distintas fases de elaboraci�n del PFC.


\subsection{Objetivo principal} \label{sec:objetivo-principal}

El PFC consiste en la elaboraci�n de una \textbf{aplicaci�n web} para la visualizaci�n de los datos obtenidos como resultado de las operaciones de medici�n y evaluaci�n de la calidad del producto software del entorno MEDUSAS, que est� basado en la nueva familia de normas ISO/IEC 25000. Para ello, se desarrollar� un proyecto Java que ser� capaz de representar gr�ficamente los distintos niveles de detalle procedentes de los resultados de evaluaci�n de la calidad del producto software.

Dicha aplicaci�n web ser� desarrollada utilizando las �ltimas tecnolog�as empleadas en desarrollo y dise�o web (Ver Secci�n \ref{sec:medios}) con el fin de obtener una RIA (Rich Internet Application) \cite{def:ria}.

Una RIA se fundamenta en una arquitectura cliente/servidor as�ncrona, segura y escalable, disminuyendo la carga de trabajo en el servidor y siendo accesible desde cualquier navegador web. As�, estar� disponible para cualquier persona desde cualquier rinc�n del planeta, tratando de maximizar la experiencia de usuario y producir un incremento en la productividad.


\subsection{Objetivos secundarios} \label{sec:objetivos-secundarios}

A lo largo del proceso de investigaci�n y elaboraci�n del PFC se alcanzar�n las siguientes metas:
\begin{myitemize}
	\item Comprender las actuales normas sobre la calidad del software.
	\item Investigar sobre las principales formas de visualizaci�n de calidad del software.
	\item Estudiar los principales entornos y medios actuales para la visualizaci�n de la calidad del software, as� como las librer�as gr�ficas que permiten su implementaci�n.
	\item Implementar un entorno que permita visualizar la calidad de un producto software de acuerdo a las caracter�sticas de calidad impuestas por la ISO/IEC 25000.
	\item Integrar dicho entorno de visualizaci�n con un entorno real de medici�n de la calidad del producto software.
\end{myitemize}
\clearpage
\section{M�todo y fases de trabajo}

Dada la naturaleza del proyecto, en el que se ir�n definiendo requisitos y funcionalidades a lo largo de todo su ciclo de vida, se ha optado por utilizar una metodolog�a de desarrollo de software gen�rica que permita adaptarse a este caso de estudio. Por estas razones, se ha seleccionado el \textbf{Proceso Unificado de Desarrollo} (en adelante PUD) como metodolog�a de trabajo.

El PUD \cite{jacobson2000proceso} es una evoluci�n del Proceso Unificado de Rational, que define un ``conjunto de actividades necesarias para transformar los requisitos de usuario en un sistema software''. Es un marco gen�rico que puede especializarse para una gran variedad de sistemas de software cualesquiera que sean el �rea de aplicaci�n o el tama�o del proyecto. Sus principales caracter�sticas son las siguientes:
\begin{myitemize}
	\item \textbf{Dirigido por casos de uso}. Para poder desarrollar un sistema es necesario saber qu� necesitan sus usuarios. Un usuario puede ser un ser humano u otro sistema que interacciona con el sistema que se est� desarrollando. Las necesidades de un usuario se denominan requisitos funcionales y se representan por medio de casos de uso.
	
	Los casos de uso gu�an el proceso de desarrollo desde la especificaci�n de requisitos hasta las pruebas y se utilizan para crear los modelos que permitan la construcci�n e implementaci�n de los mismos. Todos los casos de uso juntos constituyen el \textbf{modelo de casos de uso}.
	
	\item \textbf{Centrado en la arquitectura}. Un sistema software puede contemplarse desde varios puntos de vista. La arquitectura software incluye los aspectos est�ticos y din�micos m�s significativos del sistema y debe estar profundamente relacionada con los casos de uso ya que debe permitir el desarrollo de los mismos. Por esta raz�n, la arquitectura y los casos de uso deben desarrollarse en paralelo.
	
	\item \textbf{Iterativo e incremental}. Partiendo del dicho latino \textit{divide et impera} (en castellano, \textit{divide y vencer�s}), el ciclo de vida del proyecto (Fig. \ref{fig:mapa-conceptual-pud}) se divide en ciclos. Cada ciclo consta de cuatro fases que a su vez se dividen en iteraciones. Las fases de un ciclo son las siguientes:
	\begin{myitemize}
		\item \textbf{Inicio o Factibilidad}. Se obtiene un modelo de casos de uso simplificado, se identifican riesgos potenciales y se estima el proyecto de manera aproximada.
		
		\item \textbf{Elaboraci�n}. Se mejora el modelo de casos de uso y se dise�a la arquitectura del sistema. A continuaci�n, se desarrollan los casos de uso m�s cr�ticos que se identificaron en la fase de inicio y se obtiene una \textit{l�nea base} de la arquitectura para planificar las actividades y estimar los recursos necesarios para terminar el proyecto.
		
		\item \textbf{Construcci�n}. Se crea el producto en base a las dos fases anteriores.
		
		\item \textbf{Transici�n}. Implica la correcci�n de errores y \textit{bugs}, as� como el mantenimiento del sistema.
	\end{myitemize}
	
	 En cierto modo, una iteraci�n se corresponde con un mini-proyecto que resulta en un \textit{incremento}. En cada iteraci�n se aborda un conjunto de casos de uso que ampl�an la funcionalidad global del producto desarrollado y trata los riesgos m�s importantes. Si la iteraci�n cumple sus objetivos se contin�a con la pr�xima, en caso contrario se vuelve a las primeras etapas de la misma y se toman otras decisiones que permitan cumplir sus objetivos. Las iteraciones deben ser seleccionadas y ejecutarse de forma controlada y siguiendo un flujo de trabajo. Este flujo de trabajo consta de cinco etapas:
	\begin{myitemize}
		\item \textbf{Requisitos}. El equipo de desarrollo y los clientes establecen comunicaciones por diferentes medios para especificar qu� es lo que el usuario espera del sistema que se va a desarrollar.
		
		\item \textbf{An�lisis}. Se identifican y especifican los casos de uso relevantes.
		
		\item \textbf{Dise�o}. Se crea un dise�o utilizando la arquitectura seleccionada como gu�a para tratar de dotar al sistema de las funcionalidades representadas por los casos de uso identificados en la etapa de an�lisis.
		
		\item \textbf{Implementaci�n}. Se implementan las decisiones tomadas en la etapa de dise�o, con el fin de construir las funcionalidades que satisfacen los casos de uso identificados en la etapa de an�lisis.
		
		\item \textbf{Pruebas}. Se verifica que los cambios realizados en el sistema satisfacen los casos de uso.
	\end{myitemize}
\end{myitemize}


\imagenBorde{./images/mapa-conceptual-pud.png}{0.39}{Mapa conceptual del PUD}{fig:mapa-conceptual-pud}

\clearpage
\section{Medios que se pretenden utilizar} \label{sec:medios}

Los medios y recursos necesarios para el desarrollo y despliegue del PFC se categorizan en \textit{medios hardware} y \textit{medios software}.


\subsection{Medios hardware}

Tanto para el desarrollo como el despliegue del sistema bastar� con un ordenador sobremesa o port�til. Dicho ordenador deber� estar dotado de una conexi�n a Internet y actuar como servidor web para permitir que los usuarios puedan acceder a la aplicaci�n por medio de un navegador web convencional.

La tarjeta gr�fica del sistema requiere especial menci�n. �sta no tiene por qu� poseer una potencia extraordinaria, pero es de vital importancia que disponga de sus controladores gr�ficos debidamente instalados para una correcta ejecuci�n del motor gr�fico.


\subsection{Medios software}

Tanto la aplicaci�n web como el motor gr�fico ser�n desarrollados en tecnolog�as basadas en Java y, debido a su naturaleza \textit{multiplataforma}, se espera que no existan restricciones en cuanto al sistema operativo que el ordenador debe ejecutar. Sin embargo, �ste debe tener debidamente instalado el entorno de ejecuci�n Java (JRE o \textit{Java Runtime Environment})

Por otro lado, son necesarias las siguientes herramientas:
\begin{myitemize}
	\item JDK 1.4+ (\textit{Java Development Kit}).
	\item Un entorno integrado de desarrollo (IDE o \textit{Integrated Development Environment}). Se utilizar� \textbf{Eclipse EE Helios}.
	\item Un navegador web con JavaScript habilitado para que se pueda establecer una comunicaci�n entre la aplicaci�n web y el motor gr�fico, que estar� incrustado en un applet, y poder beneficiarse de las caracter�sticas de la tecnolog�a AJAX (\textit{Asynchronous JavaScript And XML}).
	\item Un sistema para la gesti�n y construcci�n de proyectos. Se har� uso de \textbf{Maven 2}.
	\item Un servidor web capaz de desplegar aplicaciones web con estructura WAR (\textit{Web-Archive} ). \textbf{Apache Tomcat 7} es la soluci�n m�s factible gracias a su licencia, eficiencia y rendimiento.
	\item Un sistema gestor de bases de datos. Se emplear� \textbf{MySQL Community Server 5.x}.
	\item \textbf{Visual Paradigm} para modelado y dise�o de diagramas UML.
	\item \LaTeX{} para la elaboraci�n de documentaci�n t�cnica asociada al proyecto.
	\item GIMP e Inkscape para el dise�o y edici�n de figuras, im�genes y gr�ficos.
\end{myitemize}

\clearpage

\phantomsection
%% Now actually use the newly defined style.
\urlstyle{leo}
\pagestyle{plain}
% Bibliograf�a
\selectbiblanguage{spanish}
\renewcommand{\bibname}{Bibliograf�a}
\renewcommand{\refname}{\bibname}
\addcontentsline{toc}{section}{\refname} % A�ade la bibliograf�a al �ndice
\bibliographystyle{bababbrv} % Estilo de bibliograf�a
\bibliography{Bibliografia}
\clearpage
\pagestyle{empty}
\newsavebox{\caja}
\sbox{\caja}{\textbf{Escuela Superior de Inform�tica}}
\newlength{\anchoCaja}
\setlength{\anchoCaja}{\wd\caja}
\noindent
\begin{minipage}{\anchoCaja}
	\begin{center}
	\noindent\includegraphics[keepaspectratio,scale=0.25]{./images/esi_bw.png}\\
		\noindent\textbf{Escuela Superior de Inform�tica}
	\end{center}
\end{minipage}

\vspace{10mm}
\noindent D. Jose Domingo L�pez L�pez, con D.N.I. n\textordmasculine{} 71219116-F, alumno de la Escuela Superior de Inform�tica.

\vspace{5mm}
\noindent EXPONE:\\
\textemdash{} Que se encuentra matriculado del PFC.

\vspace{5mm}
\noindent SOLICITA:\\
\textemdash{} Que la Comisi�n Acad�mica eval�e el anteproyecto de fin de carrera que acompa�a, titulado \textit{DESGLOSA: Un sistema de visualizaci�n 3D para dar soporte al Desarrollo Global de Software}.

\vspace{10mm}
\noindent Ciudad Real, a \today{}\\

\vspace{15mm}
\newlength{\factor}
\setlength{\factor}{0.33\textwidth}
\noindent 
\hfill
\begin{minipage}{\factor}
El alumno,\\
\\
\\
\\
\\
Fdo.:
\end{minipage}
\hfill
\begin{minipage}{\factor}
V\textordmasculine B\textordmasculine{} El director,\\
\\
\\
\\
\\
Fdo.:
\end{minipage}
\hfill

\vspace{40mm}
\begin{center}
	\noindent SR. DIRECTOR DE LA ESCUELA SUPERIOR DE INFORM�TICA
\end{center}
\end{document}
